\documentclass[12pt,a4paper]{article}

\usepackage{amsmath,amssymb,amsthm}
\usepackage{geometry}
\usepackage{hyperref}
\usepackage{booktabs}
\usepackage{enumitem}
\usepackage{mathtools}
\usepackage{bm}
\usepackage{xcolor}

\geometry{margin=2.2cm}

\hypersetup{
  colorlinks=true,
  linkcolor=blue!70!black,
  urlcolor=blue!60!black,
  citecolor=green!50!black
}

\newtheorem{theorem}{Theorem}[section]
\newtheorem{corollary}[theorem]{Corollary}
\newtheorem{lemma}[theorem]{Lemma}
\newtheorem{proposition}[theorem]{Proposition}
\theoremstyle{definition}
\newtheorem{definition}[theorem]{Definition}
\theoremstyle{remark}
\newtheorem{remark}[theorem]{Remark}

\newcommand{\res}{\mathcal{R}}
\newcommand{\CC}{\mathbb{C}}
\newcommand{\RR}{\mathbb{R}}
\newcommand{\ZZ}{\mathbb{Z}}
\newcommand{\HH}{\mathbf{H}}
\newcommand{\GG}{\mathbf{G}}
\newcommand{\FF}{\mathbf{F}}

\title{%
  \textbf{IQS-888: The Imaginary Extension}\\[6pt]
  \large Complex Hexagram Coordinates and the Möbius Unification\\[4pt]
  \normalsize Companion to \emph{A Unified Theory of Quantum Gravity on the 64-State Hexagram Lattice}
}
\author{Alberto Valido Delgado}
\date{March 2026}

\begin{document}
\maketitle

\begin{abstract}
The original IQS-888 paper established that quantum mechanics and general
relativity are Fourier duals of one resonance matrix $\mathbf{R}$ on the
64-state hexagram lattice. Here we show that the lattice possesses a
natural complexification: the upper and lower trigrams of each hexagram
are the real and imaginary parts of a single complex trigram. This doubles
the lattice from $\RR^{64}$ to $\CC^{64}$ and introduces the Lorentz factor
$\gamma = 1/\sqrt{1 - a^2}$ (where $a$ is the astrocyte meta-variable) as
the bridge between the real and imaginary planes. The Möbius topology of
the complex lattice eliminates the need for separate matter and antimatter
sectors, resolves the complement anomaly at Hamming distance 6, and places
the decoherence function (Luna, $e$) as the rotation operator between planes.
Every theorem in the original paper acquires an imaginary counterpart; we
show that these pairs are not independent but are the same theorem viewed
from opposite sides of a single twisted surface.
\end{abstract}

\tableofcontents

%=========================================================
\section{The Complex Trigram}\label{sec:complex-trigram}

\begin{definition}[Complex Trigram]
A \emph{complex trigram} is a 3-bit complex number
$\tau = t_{\mathrm{re}} + t_{\mathrm{im}}\,i$ where
$t_{\mathrm{re}}, t_{\mathrm{im}} \in \{0,1\}^3$. The state space of
complex trigrams has $8 \times 8 = 64$ elements.
\end{definition}

\begin{remark}
In the original paper, a hexagram is $h = 8\,t_{\mathrm{lower}} + t_{\mathrm{upper}}$
with lower and upper trigrams stacked vertically. The complex trigram
identifies $t_{\mathrm{lower}} = t_{\mathrm{re}}$ (the real part, what IS)
and $t_{\mathrm{upper}} = t_{\mathrm{im}}$ (the imaginary part, what is
BECOMING). A hexagram is not two trigrams stacked. It is one complex number.
\end{remark}

\begin{proposition}[Hexagram--Complex Trigram Isomorphism]
The map $\phi: \{0,\ldots,63\} \to \ZZ[i] / (8,8i)$ defined by
$\phi(h) = (h \bmod 8) + \lfloor h/8 \rfloor \cdot i$ is a bijection.
The Ba Gua Square is the Argand diagram of the complex trigram lattice.
\end{proposition}

\begin{proof}
Each hexagram $h \in \{0,\ldots,63\}$ decomposes uniquely as $h = 8q + r$
with $r \in \{0,\ldots,7\}$ and $q \in \{0,\ldots,7\}$. The map sends
$h \mapsto r + qi$, which is invertible: $\phi^{-1}(a+bi) = 8b + a$.
The 64 complex trigrams tile the $8 \times 8$ grid $\{0,\ldots,7\}^2$
in the Gaussian integer plane, which is exactly the Ba Gua Square.
\end{proof}

%=========================================================
\section{The Two Diagonals}\label{sec:diagonals}

In the original paper, four domains organize all elements: morph (dream),
work (build), salt (archive), vault (protect). These four domains form two
diagonal axes through a central point:

\begin{definition}[Domain Diagonals]
The \emph{dream-archive diagonal} connects morph and salt:
$d_1: \mathrm{morph} \leftrightarrow \mathrm{salt}$.
The \emph{build-seal diagonal} connects work and vault:
$d_2: \mathrm{work} \leftrightarrow \mathrm{vault}$.
The intersection of $d_1$ and $d_2$ is the \emph{origin} --- the factory,
the observer, the astrocyte at rest.
\end{definition}

\begin{proposition}[Diagonals as Imaginary Axes]
The two diagonals carry the imaginary components of all coordinates.
An element at real coordinate $\mathrm{morph}.n$ has imaginary coordinate
$\mathrm{salt}.n\,i$, and vice versa. An element at $\mathrm{work}.n$
has imaginary coordinate $\mathrm{vault}.n\,i$, and vice versa.
The full coordinate of any element is:
\begin{equation}\label{eq:complex-coord}
  z = x_{\mathrm{domain}} + y_{\mathrm{mirror}}\,i
\end{equation}
where $x$ and $y$ are positions within opposite domains.
\end{proposition}

%=========================================================
\section{The Lorentz Factor as Domain Bridge}\label{sec:lorentz}

\begin{definition}[Astrocytic Lorentz Factor]
For an element with astrocyte value $a \in [0, 1)$, the Lorentz factor is:
\begin{equation}\label{eq:gamma}
  \gamma = \frac{1}{\sqrt{1 - a^2}}.
\end{equation}
\end{definition}

\begin{theorem}[Lorentz Transformation of Coordinates]\label{thm:lorentz}
The real and imaginary coordinates of an element transform under change
of astrocyte as:
\begin{equation}\label{eq:lorentz-transform}
  \begin{pmatrix} x' \\ y' \end{pmatrix}
  = \begin{pmatrix} \gamma & -\gamma a \\ -\gamma a & \gamma \end{pmatrix}
  \begin{pmatrix} x \\ y \end{pmatrix}
\end{equation}
At $a = 0$: identity (pure real, deterministic, vault-like).
At $a \to 1$: $\gamma \to \infty$ and the element exists equally in both
planes (pure quantum superposition of domain and its mirror).
\end{theorem}

\begin{proof}
This is the standard Lorentz boost with velocity parameter $a$ (the
astrocyte) replacing $v/c$. The invariant is:
\begin{equation}
  x^2 - y^2 = x'^2 - y'^2
\end{equation}
which preserves the \emph{hyperbolic distance} between the real and
imaginary planes. The astrocyte is velocity through the imaginary plane.
At $a = 0$ the element is at rest in the real plane. At $a \to 1$ it
approaches the light cone --- the boundary where real and imaginary are
indistinguishable.
\end{proof}

\begin{remark}
In the original paper, the astrocyte was a meta-variable controlling
probability spread. Now it is revealed as the \emph{rapidity} of an element's
motion between the real and imaginary planes. The functions
\texttt{sample()}, \texttt{collapse()}, and \texttt{evolve()} in
\texttt{dodecahedron.js} are discrete Lorentz transformations.
\end{remark}

%=========================================================
\section{Correspondence Table: Original $\leftrightarrow$ Complex Extension}
\label{sec:correspondence}

For each theorem in the original paper, we state its imaginary counterpart
and show they are two views of one object.

\bigskip

%--- Theorem 1 ---
\subsection{Theorem 1: Discrete Energy Spectrum}

\noindent\textbf{Original.}
The eigenvalues of $\mathbf{R}$ (real, symmetric, $64 \times 64$)
form a discrete spectrum $\{\lambda_1, \ldots, \lambda_{64}\}$.

\medskip\noindent\textbf{Complex extension.}
The resonance matrix on the complex trigram lattice is
$\mathbf{R}_\CC = \mathbf{R} + i\,\mathbf{R}^*$, where $\mathbf{R}^*$
is the complement resonance (resonance between an element and its
imaginary mirror). The eigenvalues become:
\begin{equation}
  \lambda_n^\CC = \lambda_n + i\,\lambda_n^*
\end{equation}
The real part is the energy level. The imaginary part is the
\emph{decay rate} --- how fast the state decoherences from quantum
(imaginary) to classical (real). States with $|\lambda_n^*| \approx 0$
are stable (salt-like). States with large $|\lambda_n^*|$ are transient
(morph-like).

\medskip\noindent\textbf{Correspondence.}
$\mathrm{Re}(\lambda^\CC)$ = original spectrum. $\mathrm{Im}(\lambda^\CC)$
= decoherence spectrum. Together: the \emph{complex energy} of each state.
Luna ($e$) governs the imaginary decay: $e^{-|\lambda^*| t}$ is the
survival amplitude. The original paper's spectrum is the shadow of the
complex spectrum projected onto the real axis.

\bigskip

%--- Theorem 2 ---
\subsection{Theorem 2: Uncertainty Principle}

\noindent\textbf{Original.}
$\Delta x \cdot \Delta k \geq N / 2\pi = 32/\pi \approx 10.19$.

\medskip\noindent\textbf{Complex extension.}
On the complex lattice, position $z = x + iy$ and momentum
$\zeta = k_x + ik_y$ are both complex. The uncertainty principle
becomes:
\begin{equation}\label{eq:complex-uncertainty}
  \Delta z \cdot \Delta \zeta \geq \frac{N}{\pi}
  = \frac{64}{\pi} \approx 20.37
\end{equation}
The bound doubles because the complex plane has twice the degrees of
freedom. The imaginary part of position ($y$) is uncertain in the
imaginary part of momentum ($k_y$), independently of the real pair.

\medskip\noindent\textbf{Correspondence.}
Projecting to the real axis: $\Delta x \cdot \Delta k_x \geq 32/\pi$
(recovers the original). Projecting to the imaginary axis:
$\Delta y \cdot \Delta k_y \geq 32/\pi$ (the mirror theorem). The two
projections are independent --- you can know real position exactly while
imaginary position is maximally uncertain (a vault element: $a \approx 0$,
$\gamma \approx 1$).

\bigskip

%--- Theorem 3 ---
\subsection{Theorem 3: Wave-Particle Duality}

\noindent\textbf{Original.}
A state localized in position (Ba Gua Square) is delocalized in momentum
(King Wen Wheel), connected by the DFT $\mathbf{F}$.

\medskip\noindent\textbf{Complex extension.}
The DFT kernel $e^{-2\pi i kn/N}$ is already complex --- it rotates between
real and imaginary with each step. In the complex trigram lattice, the DFT
becomes a \emph{double DFT}:
\begin{equation}
  \hat{\psi}(k_x, k_y) = \frac{1}{N} \sum_{x,y=0}^{7}
  \psi(x + iy) \, e^{-2\pi i (k_x x + k_y y) / 8}
\end{equation}
This is a 2D Fourier transform on the $8 \times 8$ grid. The Ba Gua Square
is the position-space image. The King Wen Wheel is the
momentum-space image. They are related by rotation in the complex plane,
not just permutation.

\medskip\noindent\textbf{Correspondence.}
Wave-particle duality in the original paper operated on a 1D ring of 64
states. In the complex extension it operates on a 2D torus ($8 \times 8$
with periodic boundaries). The torus IS the Möbius surface --- the twist
comes from the fact that traversing one full cycle of the King Wen Wheel
returns you to the complement ($h \to \bar{h}$), not the identity. After
TWO cycles you return to $h$. This is the definition of a Möbius strip:
a surface where one traversal reverses orientation.

\bigskip

%--- Theorem 4 ---
\subsection{Theorem 4: Schrödinger Equation}

\noindent\textbf{Original.}
$i\hbar \, d\psi/dt = \mathbf{H}\psi$, with unitary evolution
$\mathbf{U}(t) = e^{-i\mathbf{H}t/\hbar}$.

\medskip\noindent\textbf{Complex extension.}
The Hamiltonian becomes complex-Hermitian:
$\mathbf{H}_\CC = \mathbf{H}_{\mathrm{re}} + i\,\mathbf{H}_{\mathrm{im}}$.
The real part generates oscillation (quantum mechanics). The imaginary part
generates \emph{decay} (decoherence). The evolution operator:
\begin{equation}\label{eq:complex-evolution}
  \mathbf{U}_\CC(t) = e^{-i\mathbf{H}_{\mathrm{re}} t/\hbar}
  \cdot e^{-\mathbf{H}_{\mathrm{im}} t/\hbar}
\end{equation}
The first factor is unitary (preserves probability). The second factor is
a contraction (decays probability). Together: the state oscillates AND
decays. This is the Lindblad master equation made algebraic.

\medskip\noindent\textbf{Correspondence.}
$\mathbf{H}_{\mathrm{im}} = 0$: recovers the original (closed system,
no decoherence, probability conserved forever). $\mathbf{H}_{\mathrm{im}}
\neq 0$: open system, state decays from imaginary (quantum) to real
(classical). The decay rate is $\Gamma = \|\mathbf{H}_{\mathrm{im}}\|$.
In L7: $\Gamma$ scales with the astrocyte. High astrocyte = fast decay =
rapid manifestation from dream to archive. Low astrocyte = slow decay =
stable quantum superposition preserved.

\emph{Luna is} $e^{-\Gamma t}$: \emph{the decoherence function in the
imaginary plane is the survival probability in the real plane.}

\bigskip

%--- Theorem 5 ---
\subsection{Theorem 5: Inverse-Square Gravity with Resonance Correction}

\noindent\textbf{Original.}
$F = G m_1 m_2 \res(h_1, h_2) / r^2$.

\medskip\noindent\textbf{Complex extension.}
The distance $r$ between two hexagrams becomes complex:
\begin{equation}
  r_\CC = d_{\mathrm{Hamming}}(h_1, h_2) +
  i \cdot d_{\mathrm{Hamming}}(\bar{h}_1, \bar{h}_2)
\end{equation}
where $\bar{h}$ is the complement (imaginary mirror). The gravitational
force:
\begin{equation}\label{eq:complex-gravity}
  F_\CC = \frac{G \, m_1 \, m_2 \, \res_\CC(h_1, h_2)}{|r_\CC|^2}
\end{equation}
where $|r_\CC|^2 = r_{\mathrm{re}}^2 + r_{\mathrm{im}}^2$.

\medskip\noindent\textbf{Correspondence.}
The complement anomaly at Hamming distance 6 is \emph{resolved}: $h$ and
$\bar{h}$ are at maximum real distance ($d=6$) but minimum imaginary
distance ($d_{\mathrm{im}} = 0$, because the complement of a complement is
the original). The enhanced resonance $\gamma = 0.12$ at $d=6$ is not
anomalous --- it's the imaginary distance being zero. The wormhole
(ER\,=\,EPR) is a path through the imaginary plane that bypasses real
distance entirely.

\bigskip

%--- Theorem 6 ---
\subsection{Theorem 6: Mass-Energy Equivalence}

\noindent\textbf{Original.}
$E = mc^2 = 4096\,m$, with $c = 64$ and $c^2 = 4096 = 4^6 = 2^{12}$.

\medskip\noindent\textbf{Complex extension.}
On the complex lattice, the speed of light becomes complex:
\begin{equation}
  c_\CC = c + ic = 64(1 + i), \qquad |c_\CC|^2 = 2c^2 = 8192.
\end{equation}
Mass-energy equivalence:
\begin{equation}\label{eq:complex-emc2}
  E_\CC = m \cdot |c_\CC|^2 = 8192\,m.
\end{equation}
This is \emph{exactly} the pair annihilation energy from the original paper:
$E_{\mathrm{pair}} = 2mc^2 = 8192\,m$.

\medskip\noindent\textbf{Correspondence.}
In the original paper, pair annihilation of $(h, \bar{h})$ released
$2mc^2$. In the complex extension, a \emph{single} particle on the complex
lattice already carries $|c_\CC|^2 = 2c^2$ because it exists in both the
real and imaginary planes simultaneously. Pair annihilation is not two
particles destroying each other --- it's one complex particle collapsing
from $\CC$ to $\RR$, losing its imaginary component. The factor of 2 was
always the imaginary axis.

\bigskip

%--- Theorem 7 ---
\subsection{Theorem 7: Discrete Einstein Field Equation}

\noindent\textbf{Original.}
$\mathcal{R}(h) = \kappa \cdot T(h)$, with curvature monotonically
decreasing with Hamming distance (except the complement anomaly).

\medskip\noindent\textbf{Complex extension.}
The curvature tensor is complex:
\begin{equation}
  \mathcal{R}_\CC(z) = \mathcal{R}_{\mathrm{re}}(x) +
  i\,\mathcal{R}_{\mathrm{im}}(y)
\end{equation}
Real curvature bends space (gravity). Imaginary curvature bends
\emph{time} (decoherence). The Einstein equation becomes:
\begin{equation}\label{eq:complex-einstein}
  \mathcal{R}_\CC(z) = \kappa \cdot T_\CC(z)
\end{equation}
where $T_\CC = T_{\mathrm{mass}} + i\,T_{\mathrm{entropy}}$. Mass curves
space; entropy curves time. This is the first equation in which
\emph{entropy enters general relativity as the imaginary part of the
stress-energy tensor}.

\medskip\noindent\textbf{Correspondence.}
At $a = 0$ (vault): $\mathcal{R}_{\mathrm{im}} = 0$, pure spatial
curvature, classical GR recovered. At $a \to 1$ (morph):
$\mathcal{R}_{\mathrm{im}}$ dominates, spacetime curvature is primarily
temporal, the element is ``falling through time'' --- decoherencing from
quantum to classical. The Schwarzschild radius in the complex extension
has an imaginary part: the event horizon is not a surface in space but a
\emph{cone in spacetime} --- the hourglass.

\bigskip

%--- Theorem 8 (Main) ---
\subsection{Theorem 8: Unified Quantum Gravity (The Emerald Equation)}

\noindent\textbf{Original.}
$\mathbf{H}_{\mathrm{quantum}} = \mathbf{F} \cdot \mathbf{G}_{\mathrm{gravity}} \cdot \mathbf{F}^{-1}$.

\medskip\noindent\textbf{Complex extension.}

\begin{theorem}[The Emerald Equation]\label{thm:emerald}
On the complex hexagram lattice $\CC^{64}$, quantum mechanics, general
relativity, and decoherence are three projections of one complex resonance
matrix:
\begin{equation}\label{eq:emerald}
  \boxed{%
    \mathbf{H}_\CC = \mathbf{F}_\CC \cdot \mathbf{G}_\CC
    \cdot \mathbf{F}_\CC^{-1}
    = (\mathbf{H} + i\Gamma) = \mathbf{F}(\mathbf{G} + i\mathbf{S})\mathbf{F}^{-1}
  }
\end{equation}
where:
\begin{itemize}[nosep]
  \item $\mathbf{G}$ = gravitational coupling (real part, position basis)
  \item $\mathbf{S}$ = entropic coupling (imaginary part, position basis)
  \item $\mathbf{H}$ = quantum Hamiltonian (real part, momentum basis)
  \item $\Gamma$ = decoherence operator (imaginary part, momentum basis)
  \item $\mathbf{F}_\CC$ = 2D DFT on the $8 \times 8$ complex trigram lattice
\end{itemize}
\end{theorem}

\begin{proof}
The original proof showed $\mathrm{spec}(\mathbf{H}) = \mathrm{spec}(\mathbf{G})$
because the DFT is unitary. Extend $\mathbf{R}$ to $\mathbf{R}_\CC =
\mathbf{R} + i\mathbf{R}^*$ where $\mathbf{R}^*_{mn} = \res(\bar{m}, \bar{n})$
is the complement resonance matrix.

$\mathbf{R}^*$ is also real and symmetric (complement is an involution), so
$\mathbf{R}_\CC$ is complex-Hermitian:
$\mathbf{R}_\CC^\dagger = (\mathbf{R} + i\mathbf{R}^*)^\dagger
= \mathbf{R} - i\mathbf{R}^* + 2i\mathbf{R}^* = \mathbf{R}_\CC$
(since $\mathbf{R}$ and $\mathbf{R}^*$ are both real symmetric).

The 2D DFT $\mathbf{F}_\CC$ on $\CC^{64}$ is unitary (product of two unitary
1D DFTs). Therefore:
\begin{equation}
  \mathrm{spec}(\mathbf{F}_\CC \mathbf{R}_\CC \mathbf{F}_\CC^{-1})
  = \mathrm{spec}(\mathbf{R}_\CC).
\end{equation}
The complex eigenvalues $\lambda_n + i\mu_n$ encode both the energy ($\lambda_n$)
and the decoherence rate ($\mu_n$) of each state. The three sectors ---
quantum mechanics (momentum basis, real part), general relativity (position
basis, real part), and decoherence (imaginary part in either basis) --- share
one complex spectrum.

As above, so below. One matrix. One spectrum. Three shadows.
\end{proof}

%=========================================================
\section{The Four Corollaries Extended}\label{sec:corollaries-ext}

\begin{corollary}[No Singularities --- Complex]\label{cor:sing-complex}
On the complex lattice, the minimum volume is one complex cell: $1 + i$.
The Planck length acquires an imaginary part: $\ell_P^\CC = \ell_P(1+i)$.
Black hole singularities cannot form in either the real or imaginary
plane --- information is preserved as a complex hexagram bit pattern.
\end{corollary}

\begin{corollary}[ER\,=\,EPR --- Complex]\label{cor:erepr-complex}
The complement bond $(h, \bar{h})$ is a path of zero imaginary distance.
The wormhole IS the imaginary axis: two points at maximum real separation
are adjacent in the imaginary plane. This is not tunneling --- it is
navigation. The Möbius twist maps $h$ to $\bar{h}$ in one half-traversal
of the surface.
\end{corollary}

\begin{corollary}[Matter-Antimatter --- Complex]\label{cor:matter-complex}
Matter ($h$) and antimatter ($\bar{h}$) are the real and imaginary
projections of one complex hexagram. They are not two particles but one
particle viewed from two sides of the Möbius surface. Pair annihilation
= collapse of the imaginary axis. The released energy $8192\,m = 2mc^2
= m|c_\CC|^2$ is the full complex energy of one particle, not the
combined energy of two.
\end{corollary}

\begin{corollary}[12 $\to$ 24 Dimensions]\label{cor:24d}
Classical hexagrams: $2^6 = 64$ states (6 real dimensions).
Quantum hexagrams: $4^6 = 4096$ states (12 dimensions).
Complex quantum hexagrams: $4096 + 4096i = 4096(1+i)$ states
(12 real + 12 imaginary = 24 dimensions).

$24 = $ the dimension of the Leech lattice, the densest sphere packing
in 24 dimensions, the lattice underlying the Monster group, and the
number of dimensions in bosonic string theory. The complex hexagram
lattice IS the Leech lattice reduced to its binary skeleton.
\end{corollary}

%=========================================================
\section{The Number 666}\label{sec:666}

Every element in the lattice is addressed by three coordinates:
\begin{equation}
  z = x_{\mathrm{real}} + y_{\mathrm{imaginary}} \cdot i,
  \quad \gamma = \frac{1}{\sqrt{1 - a^2}}
\end{equation}

The Sun magic square is $6 \times 6$. Each row sums to 111. The total
is 666. Three sixes: real (6 bits), imaginary (6 bits), gamma (the 6th
constant, the bridge). The hexagram is 6 lines. Carbon --- the atom of
life --- has 6 protons, 6 neutrons, 6 electrons. The Star of David is
two interlocking triangles (real and imaginary) forming a hexagram.

\begin{equation}
  666 = \sum_{n=1}^{6} \binom{6}{n} \cdot n \cdot (7-n)
  = 6 + 30 + 60 + 60 + 30 + 6 + \ldots
\end{equation}

But more precisely: $6 + 6 + 6 = 18$ bits. The Q64 encoding already uses
18 bits of contextual addressing (6 I-Ching + 8 Ifa + 4 Geomancy). The
three sixes are not Satan. They are carbon. They are life. They are the
address of every atom in creation.

%=========================================================
\section{Luna Completes the Constants}\label{sec:luna}

The four fundamental constants of IQS-888:
\begin{align}
  c &= 64 \quad \text{(lattice propagation speed)} \\
  \hbar &= 1 \quad \text{(minimum action quantum)} \\
  G &= 4\pi^2 \quad \text{(gravitational coupling)} \\
  e &= 2.71828\ldots \quad \text{(decoherence constant --- Luna)}
\end{align}

With the imaginary extension, their roles clarify:
\begin{itemize}[nosep]
  \item $c$ governs the \emph{real} plane (spatial propagation)
  \item $e$ governs the \emph{imaginary} plane (temporal decay / rotation)
  \item $\hbar$ is the minimum unit in \emph{both} planes
  \item $G$ couples the planes (gravity = real curvature induced by
    imaginary mass-energy)
\end{itemize}

Euler's identity ties them:
\begin{equation}
  e^{i\pi} + 1 = 0
\end{equation}
All five fundamental objects ($e$, $i$, $\pi$, $1$, $0$) appear. On the
hexagram lattice, this becomes:
\begin{equation}
  e^{i \cdot \pi \cdot \hbar \cdot c / G}
  = e^{i \cdot \pi \cdot 1 \cdot 64 / 4\pi^2}
  = e^{i \cdot 16/\pi}
  = e^{i \cdot 5.093\ldots}
\end{equation}
The phase $16/\pi \approx 5.093$ is within $0.093$ of $5$ --- the emerald
stone, the center of the Lo Shu magic square, the 5-degree shift per cycle
of the 365-day spiral.

The universe almost closes ($e^{2\pi i} = 1$) but not quite
($e^{2 \times 365\pi i / 360} \neq 1$). The remainder is the 5-degree
emerald residue. This is why time is a spiral, not a circle. This is
why the salted elements fold onto themselves but never quite reach zero.
This is why $\gamma$ is always $> 1$ when $a > 0$.

The lapis philosophorum does not create gold. It shows you the gold was
always in the lead --- the imaginary part of the real.

\end{document}
